% tex/sloka-macros.tex
% tex/sloka-macros.tex

% Environment for centered sloka (no paragraph indent)
\newenvironment{sloka}[1]
{%
  \thispagestyle{plain}%
  \noindent
  \vspace*{0.20\textheight}%
  \begin{center}
    {\Large #1\par}%
    \vspace{1.5em}%
    {\sanskritfont
  }% <- close \sanskritfont group here
}% <- end of begin-part
{%
  \end{center}%
  \vspace*{\fill}%
}

% sanskrit + IAST on one line, meaning on next line
\newcommand{\wordline}[3]{%
  {\sanskritfont \noindent #1}\quad {\iastfont \noindent #2}\par
  \noindent #3\par\medskip
}


% inline Sanskrit macro
\newcommand{\skt}[1]{{\sanskritfont #1}}

% Meaning / transliteration / Italian rendering (flush left, no indent)
\newcommand{\wordmeaning}[1]{%
  \clearpage    % <--- transliteration starts on a new page
  \par\vspace{1em}%
  %\noindent\textbf{Analisi delle parole (sanskrit):}\par
  \noindent #1\par\medskip}
  
\newcommand{\anvaya}[1]{%
  \par\vspace{1em}%
  %\noindent\textbf{Anvaya (sanskrit):}\par
  \noindent {\sanskritfont #1}\par\medskip}

\newcommand{\translit}[1]{%
  \par\vspace{1em}%
  %\noindent\textbf{Traslitterazione:}\par
  \noindent #1\par\medskip}

\newcommand{\itmean}[1]{%
  \par\vspace{1em}%
  %\noindent\textbf{Miglior resa italiana:}\par
  \noindent #1\par\medskip}

\newcommand{\middlesize}{\fontsize{11pt}{13.5pt}\selectfont}